\metatitle{Real-rooted polynomials: graphs and combinatorics}

\metadescription{Matching polynomials, rook polynomials, independence polynomials of claw-free graphs, Touchard polynomials, and Chow polynomials.}


\section[realRootedMatching]{Matching and independence polynomials}

Let $G = (V,E)$ be any graph.
The \defin{matching polynomial} $\mu_G(x)$ is defined as
\[
\mu_G(x) \coloneqq \sum_{\substack{M \subseteq E(G) \\ \text{$M$ matching}}} x^{|M|}
\]
where $|M|$ is the number of edges in the matching.

In \cite[Thm. 4.2]{HeilmannLieb1972}, it was shown that $\mu_G(x)$
is real-rooted. Moreover, $\mu_G(x)$ interlaces $\mu_{G'}(x)$, where $G'$
is any graph obtained from $G$ by removing a vertex.
This is a very powerful theorem, but it follows
from the more general Chudnovsky--Seymour theorem on independence polynomials below.

A hypergraph generalization (including multigraphs)
is considered by \name{Nima Amini} in \cite{Amini2019x}.


\begin{example*}[Laplacian matching polynomial]
Let $d_G(v)$ denote the degree of $v$ in $G$.
The \defin{Laplacian matching polynomial} is the matching-theoretic Laplacian analogue
\[
 \operatorname{LM}(G,x) \coloneqq
 \sum_{\substack{M \subseteq E(G) \\ \text{$M$ matching}}}
 (-1)^{|M|}
 \prod_{v\in V(G)\setminus V(M)} (x-d_G(v)).
\]
It is real-rooted with nonnegative roots, see \cite{WanWangMohammadian2022}.
One way to see why this belongs with matching polynomials is the subdivision identity
\[
 \mathsf{M}(S(G),x)=x^{|E(G)|-|V(G)|}\operatorname{LM}(G,x^2),
\]
where $S(G)$ is the subdivision graph and
\[
 \mathsf{M}(H,x) \coloneqq
 \sum_{\substack{M \subseteq E(H) \\ \text{$M$ matching}}}
 (-1)^{|M|}x^{|V(H)\setminus V(M)|}
\]
is the usual Heilmann--Lieb normalization of the matching polynomial.

For connected $G$, \name{Jiang-Chao Wan}, \name{Yi Wang} and \name{Ali Mohammadian} also show that $0$ is a root of
$\operatorname{LM}(G,x)$ if and only if $G$ is a tree; that the number of distinct
positive roots is at least the length of the longest path in $G$; and that
$\operatorname{LM}(G,x)$ and $\operatorname{LM}(G-e,x)$ have interlacing roots for each
edge $e$ of $G$.
\end{example*}


\begin{example*}[Fibonacci tilings via matchings]
Recall that the number of ways to tile the  one-row Young diagram $(n)$
with tiles of shape $(1)$ and $(2)$ is given by the $(n+1)th$ Fibonacci number.
If we then define
\[
 F_0(x) \coloneqq 1 \quad F_1(x) \coloneqq 1, \quad F_n(x) \coloneqq F_{n-1}(x) + x\cdot F_{n-2}(x)
\]
we see that $[x^k] F_n(x)$ is the number of ways to tile $(n)$ using exactly $k$ (2)-tiles.

By the Heilmann--Lieb theorem, we can see that for all $n$, $F_n(x)$ is real-rooted,
since we can see that $\mu_G(x) = F_n(x)$ for the path graph with $n$ vertices.
\end{example*}




\begin{example*}[Touchard polynomials (blocks in set partitions)]\label{touchardAsMatching}

The \hyperref[touchardPolynomial]{Touchard polynomials} may be defined via
\[
T_n(x) \coloneqq \sum_{P \in SP(n)} x^{blocks(P)} = \sum_{k=1}^n S(n,k) x^k.
\]
\emph{The matching approach:} \cite[Lem. 5.1]{Godsil2017}, states
that the Stirling number $S(n,k)$ is equal to the number
of size $n-k$ matchings in the bipartite graph $G_n$ on vertex set
\[
 \{1,2,\dotsc,n\} \cup \{1',2',\dotsc,n'\}
\]
where $i$ has an edge with $j'$ if and only if $i \lt j$.

For example, the graph $G_4 = \{ \{1,2'\}, \{1,3'\}, \{1,4'\}, \{2,3'\}, \{2,4'\}, \{3,4'\} \}$
has $7=S(4,4-2)$ matchings of size $2$,
and $6=S(4,4-1)$ matchings of size $1$.
We then have that $x^n T_n(x^{-1}) = \mu_{G_n}(x)$ where
$\mu_{G_n}(x)$ is the matching polynomial of $G_n$.
Since removal of vertices gives interlacing, we have that
$\mu_{G_{n-1}}(x) \interl \mu_{G_n}(x)$
from which interlacing for Touchard polynomials follows.


\emph{The rook placement approach:}
The matching polynomial $\mu_{G_n}(x)$ for the bipartite graph above, is the rook polynomial
of the staircase board of size $n-1$.

\end{example*}


\begin{example*}[r-Lah number polynomials and $r$-Touchard polynomials]

An interpretation of $r$-Lah polynomials and $r$-Stirling numbers as
matching polynomials is given in \cite{NyulRacz2021}.
\end{example*}



\subsection[rookPolynomials]{Rook polynomials}

The \defin{rook polynomial} of a board $B$ is defined as
\[
  R_B(z) \coloneqq \sum_k r_k(B) z^k
\]
where $r_k(B)$ is the number of ways to place $k$ non-attacking rooks on the board.
Here, the board can be any subset of some $n {\times} n$-board.
Note that $R_B(x)$ is exactly a matching polynomial, $\mu_G(x)$ for a bipartite
graph $G$ on vertex set $\{v_1,\dotsc,v_n,v'_1,\dotsc,v'_n\}$
where  $(v_i,v'_j)$ is an edge if and only if $(i,j)$ is a cell of the board $B$.

Nijenhuis \cite{Nijenhuis1976} proved that the rook polynomials are always real-rooted.
In fact, Nijenhuis shows a weighted result involving permanents.
\begin{theorem}[Nijenhuis, 1976]
Let $A = (a_{ij})$ be a non-negative $m\times n$-matrix.
Then
\[
 R_{A}(z) \coloneqq \sum_{\rho} (-z)^{\length(\rho)} \prod_{j=1}^{\length(\rho)}
  a(\rho_j)
\]
where we sum over all ways to place rooks on the $m\times n$-board,
and $a(\rho_j)$ is the value of the matrix at position given by the $j$th rook.
\end{theorem}


The \defin{hit polynomial} of a board can also be defined in terms of the $r_k(B)$.
Let $B$ be a subset of the $n\times n$-board.
The hit polynomial is then defined as
\[
 T_B(z) \coloneqq \sum_{k =0}^n k! r_{n-k}(B)(z-1)^{n-k} = \sum_{j=0}^n h_j(B) z^j.
\]
Here, $h_j(B)$ is the number of ways to place $n$ rooks on the $n \times n$-board,
with exactly $j$ rooks in $B$, see \cite{KaplanskyRiordan1946,Haglund2000}.

A theorem by Haglund--Ono--Wagner \cite{HaglundOnoWagner1999}
states that if $B$ is a Ferrers board (same as a Young diagram),
then $T_B(z)$ is real-rooted.




\subsection[realRootedIndependence]{Independence polynomials for claw-free graphs}

Let $G = (V,E)$ be any graph. The \defin{independence polynomial} $I(G,x)$ is defined
as
\[
I(G,x) \coloneqq \sum_{\substack{A \subseteq V(G) \\ \text{$A$ independent}}} x^{|A|}
\]
where $|A|$ is the number of edges in the matching.

A graph called \defin{claw-free} if no induced subgraph is isomorphic
to the graph with edges $\{\{1,2\},\{1,3\},\{1,4\}\}$.

If $G$ is any simple graph, its \emph{line graph} $L(G)$ is also simple.
Moreover, $L(G)$  is claw-free. Finally, we have that $I(L(G),x) = \mu_G(x)$.

\begin{theorem}[Chudnovsky--Seymour 2007, \cite{ChudnovskySeymour2007},]
Suppose $G$ is claw-free. Then $I(G,x)$ is real-rooted.
As a corollary, we get the Heilmann--Lieb theorem on matching polynomials.
\end{theorem}
\begin{proof*}
This very briefly illustrates the ideas. Note that if $v \in V(G)$ and $nbh(v)$
is $v$ and all vertices adjacent to $v$, then
\[
 I(G,x) = I(G \setminus \{v\} ,x) + x \cdot  I(G \setminus nbh(v),x).
\]
Essentially, the above sum is broken down into the cases where $v \notin A$, and when $v \in A$.
Observe that if $G$ is claw-free, then any induced subgraph is also claw-free.
Chudnovsky--Seymour then use the notion of \hyperref[compatiblePolynomials]{compatible polynomials}
to obtain their result.
\end{proof*}

This result was conjectured by \name{Richard Stanley} in \cite{Stanley1998},
a vertex-weighted version is given by \name{Alexander Engström} \cite{Engstrom2007}
and an alternative (non-vertex-deletion-recursive)
proof is given by \name{Bodo Lass} \cite{Lass2012}.
Finally, a multivariate refinement is given by \cite{LeakeRyder2019} where
one obtains \hyperref[stablePolynomial]{stable polynomials}.

For general graphs, independence polynomials need not be real-rooted.
Nevertheless, there are useful partial results.

\begin{theorem}[Brown--Dilcher--Nowakowski, \cite{BrownDilcherNowakowski2000}]
For any graph $G$, a zero of $I(G,x)$ of smallest modulus is real.
Moreover, if $G$ is well-covered with independence number $\beta(G)$,
then all zeros of $I(G,x)$ lie in the annulus
\[
  \frac{1}{|V(G)|}\leq |z|\leq \beta(G),
\]
with equality on the boundary only when $G$ is complete.
Finally, every well-covered graph $G$ is an induced subgraph of a
well-covered graph $H$ with $\beta(H)=\beta(G)$ such that $I(H,x)$ has
only simple real zeros.
\end{theorem}

The last statement shows that real-rooted independence polynomials are
plentiful among well-covered graph expansions, even though well-covered
graphs themselves may have non-real zeros. For caterpillar trees, there
are also broad unimodality results for independence polynomials, see
\cite{BahlsEthridgeSzabo2018}; these are unimodality results rather than
real-rootedness theorems.

\name{Alexander Engström} proved the following weighted generalization of the Chudnovsky--Seymour theorem:
\begin{theorem*}[A. Engström, \cite[Thm. 2.5]{Engstrom2007}]
Let $G = (V,E)$ be a graph, and $w : V \to \setR$ be vertex weights.
The \defin{weighted independence polynomial} $I(G,w,x)$ is defined as
\[
I(G,w,x) \coloneqq \sum_{\substack{A \subseteq V(G) \\ \text{$A$ independent}}} x^{|A|} \prod_{v\in A} w(v).
\]
If $G$ is claw-free and all weights are non-negative, then $I(G,w,x)$ is real-rooted.
\end{theorem*}


\subsection[realRootsChow]{Chow polynomials}

In \cite{BrandenVecchi2025x}, the authors show various
real-rootedness results regarding Chow polynomials.
